Este capítulo analisa os resultados obtidos à luz da literatura existente, interpreta suas implicações práticas, examina as limitações do estudo e contextualiza o desempenho dos modelos em relação aos objetivos propostos na Introdução.

\section{Interpretação dos Resultados Preditivos}

Os resultados apresentados no Capítulo~\ref{4_Resultados} revelaram um achado relevante sobre a comparabilidade entre paradigmas de aprendizado de máquina. Sob protocolo unificado de previsão direta $h=1$, foram empregadas as mesmas variáveis por instante, o mesmo corte da data de origem e as mesmas datas-alvo, mas com representação tabular para o XGBoost e sequencial para as redes. Nesse protocolo, XGBoost, Bi-LSTM e Bi-GRU apresentam MAEs descritivamente próximos (61,47, 61,74 e 62,10). A liderança depende da métrica: XGBoost tem o menor MAE, Bi-LSTM o menor MAPE e Bi-GRU o menor RMSE e o maior $R^2$ (0,442). Os contrastes pareados de \textit{block bootstrap}, contudo, não separam os MAEs dos três modelos; a diferença pontual não equivale a evidência de superioridade.

O resultado não sustenta a hipótese de superioridade sistemática das arquiteturas de \textit{Deep Learning} sobre modelos baseados em árvores neste conjunto de teste. A proximidade descritiva entre XGBoost e Bi-LSTM pode estar associada a três fatores complementares:

\begin{enumerate}
    \item \textbf{Engenharia de atributos temporais:} A construção manual de 40 \textit{features}, incluindo defasagens de 1, 2, 3 e 7 dias, médias móveis exponenciais e codificações harmônicas, forneceu ao XGBoost informação temporal explícita. A ablação mostrou, porém, que essa contribuição não é uniforme: histórico + calendário superou o conjunto completo no teste pós-hoc, enquanto o grupo climático não apresentou ganho incremental.

    \item \textbf{Regularização e seleção temporal:} O XGBoost dispõe de mecanismos de regularização nativos ($L_1$, $L_2$, Gamma, \texttt{max\_depth}, \texttt{subsample}) e recebeu \textit{Grid Search} com validação temporal. As redes usaram Dropout e \textit{Weight Decay}, mas não receberam busca equivalente; essa assimetria limita atribuir as diferenças apenas às arquiteturas.

    \item \textbf{Tamanho limitado do conjunto de dados:} Com 2.700 linhas de origem de treinamento (XGBoost) e 2.687 sequências (RNNs) com 41 variáveis por instante, o volume de dados é modesto para redes recorrentes profundas. O XGBoost, por sua natureza \textit{ensemble} de árvores, é adequado a dados tabulares desse porte \cite{chen2016xgboost}.
\end{enumerate}

\section{Incerteza Estatística e Relevância Prática}\label{sec:discussao_incerteza}

Os intervalos pareados apresentados na Seção~\ref{sec:resultados_bootstrap} alteram a interpretação das diferenças pontuais. O fato de um modelo registrar MAE numericamente menor não basta para afirmar que sua vantagem se manteria sob outra realização temporal do mesmo processo. Nos contrastes entre XGBoost, Bi-LSTM e Bi-GRU, todos os intervalos de 95\% contêm zero. Assim, o experimento não oferece evidência suficiente para ordenar os três modelos pelo MAE, embora permita descrever os valores observados no período de teste.

Essa conclusão também não demonstra que os modelos sejam equivalentes. Ausência de evidência de diferença e evidência de equivalência são afirmações distintas. Para sustentar equivalência prática seria necessário definir, antes da análise, uma margem de diferença considerada operacionalmente irrelevante e aplicar um procedimento inferencial formulado para essa hipótese. Como tal margem depende do custo de mobilizar equipes, do custo de um alerta não atendido e da capacidade diária da concessionária, ela não pode ser deduzida apenas das métricas estatísticas disponíveis.

Os três modelos treinados apresentaram ganhos pontuais de aproximadamente 6,5 a 7,1 interrupções por dia em relação à persistência. Entretanto, os intervalos dos contrastes semanais incluem zero e sua classificação varia nas análises de sensibilidade. Portanto, os valores centrais favorecem os modelos treinados, mas o estudo não deve converter essa vantagem observada em uma garantia geral de superioridade sobre a referência ingênua. Uma confirmação prospectiva, com previsões registradas antes da ocorrência dos alvos, reduziria a ambiguidade entre ganho circunstancial e ganho reproduzível.

A sensibilidade ao comprimento do bloco possui significado substantivo. Blocos de sete, 14 e 30 dias preservam escalas diferentes de dependência e de concentração dos erros; por isso, mudanças nos limites dos intervalos não constituem uma falha do procedimento, mas revelam quanto a conclusão depende da estrutura temporal assumida. O \textit{block bootstrap} reduz o problema de tratar dias vizinhos como independentes, porém continua condicionado ao único período observado e às escolhas de reamostragem \cite{lahiri2003resampling}.

Quando os resultados não permitem separar os modelos com segurança, outros critérios ganham importância na escolha operacional. O menor custo de treinamento, a ausência de dependência de GPU e a auditabilidade relativa tornam o XGBoost uma alternativa parcimoniosa para $h=1$. Isso não o transforma no melhor modelo em sentido universal: a Bi-GRU apresentou os menores MAEs observados em $h\geq3$, e a liderança mensal variou. A decisão deve combinar antecedência, função de custo, estabilidade prospectiva, recursos computacionais e capacidade de explicar o resultado aos responsáveis pela operação.

\section{Trade-offs entre XGBoost e Deep Learning}

A comparação entre os paradigmas vai além das métricas numéricas. A Tabela~\ref{tab:tradeoffs} sintetiza os \textit{trade-offs} práticos relevantes para uma eventual implantação operacional.

\begin{table}[H]
\caption{Trade-offs Práticos entre XGBoost e Redes Neurais Recorrentes}
\label{tab:tradeoffs}
\begin{center}
\small
\begin{tabular}{p{4cm} p{4.5cm} p{5cm}}
\toprule
\textbf{Critério} & \textbf{XGBoost} & \textbf{Bi-LSTM / Bi-GRU} \\
\midrule
Acurácia Geral ($h=1$, 365 dias) & $R^2 = 0{,}409$, MAE $= 61{,}47$ & Bi-LSTM: MAE $=61{,}74$; Bi-GRU: $R^2=0{,}442$ \\
Previsão $h>1$ (direto) & Menor MAE apenas em $h=1$ & \textbf{Bi-GRU} tem menor MAE em $h=3$, $h=7$ e $h=14$ \\
Velocidade de Treino & $<2$ min & 20--25 min (150 épocas) \\
Interpretabilidade & Alta (\textit{Feature Importance}) & Baixa (caixa-preta) \\
Necessidade de GPU & Não & Recomendável \\
Escalabilidade & Alta (paralelismo nativo) & Limitada pelo BPTT \\
Sensibilidade a Outliers & Conservador (regularização) & Maior volatilidade \\
Representação de entrada & 41 variáveis em uma linha & Mesmas 41 variáveis por instante + janela de 14 dias \\
\bottomrule
\end{tabular}
\normalsize
\end{center}
\small{Fonte: Elaborado pelos autores (2026).}
\end{table}

Uma vantagem frequentemente subestimada do XGBoost é a \textbf{auditabilidade de sua estrutura}. Conforme evidenciado pela Figura~\ref{fig:importancia} (Capítulo~\ref{4_Resultados}), o modelo fornece a importância por ganho das variáveis usadas nas árvores. Essa medida ajuda a detectar dependência excessiva de atributos, mas não identifica causalidade nem demonstra valor incremental de um grupo. Essa transparência relativa é útil em contextos regulados, desde que acompanhada por testes de ablação e validação fora da amostra \cite{rudin2019stop}.

Em contrapartida, as redes recorrentes demonstraram uma propriedade relevante na distribuição residual (Seção~\ref{sec:residuais}): suas distribuições KDE apresentam maior densidade próxima de zero. Essa forma, isoladamente, não implica menor MAE em todas as faixas; a análise segmentada mostra que o XGBoost obteve o menor MAE no baixo volume, a Bi-LSTM no volume intermediário e a Bi-GRU no alto volume. Com dados adicionais e em cenários de maior complexidade temporal (resolução horária, por exemplo), permanece como hipótese de trabalho avaliar como essas diferenças se generalizam.

\section[Inversão da Hierarquia por Horizonte]{Inversão da Hierarquia no Contexto Multi-\hspace{0pt}Horizonte}\label{sec:discussao_multihorizonte}

A análise multi-horizonte (Seção~\ref{sec:multihorizonte}, Capítulo~\ref{4_Resultados}), conduzida com o protocolo de \textbf{previsão direta} sobre 352 datas-alvo idênticas por horizonte, revelou uma nuance fundamental: \textbf{a ordenação dos MAEs observados varia com o horizonte de previsão}.

No recorte comum, em $h = 1$, os três modelos treinados apresentam MAEs próximos (62,86 para XGBoost, 63,03 para Bi-LSTM e 63,45 para Bi-GRU) e inferiores à persistência (69,83). A Bi-GRU, porém, obtém o menor RMSE (98,48) e o maior $R^2$ (0,430). Para $h = 3$, $h=7$ e $h=14$, a Bi-GRU registra o menor MAE entre os modelos treinados (70,72, 68,17 e 67,50), enquanto a Bi-LSTM tem o menor RMSE em $h=3$ e $h=7$. Todos os modelos treinados permanecem abaixo da persistência em MAE nos quatro horizontes.

\begin{itemize}
    \item O menor RMSE da \textbf{Bi-LSTM} em $h=3$ e $h=7$ indica menor penalização quadrática nesses recortes, embora a Bi-GRU tenha menor MAE; a escolha depende da função de custo operacional.
    \item Os menores MAEs da \textbf{Bi-GRU} em $h=3$, $h=7$ e $h=14$ motivam a hipótese de que sua parametrização mais compacta tenha favorecido a generalização nesse recorte; uma ablação específica seria necessária para testar essa explicação.
    \item O \textbf{XGBoost} utiliza defasagens explícitas em uma única linha, enquanto as redes recebem 14 instantes. Essa diferença de representação é parte do protocolo implementado e impede atribuir a ordenação observada exclusivamente à presença ou ausência de memória recorrente.
\end{itemize}

No conjunto avaliado, os três modelos ficaram próximos em $h=1$, enquanto a Bi-GRU registrou o menor MAE entre os modelos treinados em cada horizonte $h \geq 3$. Uma eventual combinação operacional, escolhida conforme MAE, RMSE ou MAPE e o horizonte, deve ser validada prospectivamente antes da implantação.

A decomposição mensal (Seção~\ref{sec:multihorizonte_mensal}) qualifica essa recomendação. A liderança não é uniforme ao longo do período de teste: em $h=1$, o XGBoost obtém o menor MAE em 6 dos 12 meses, a Bi-GRU em 4 e a Bi-LSTM em 2. Em $h=7$, a Bi-LSTM vence 5 meses, embora a Bi-GRU lidere no agregado. A vantagem agregada, portanto, não representa superioridade universal em todos os meses. Uma política operacional baseada apenas na média global ocultaria variações sazonais relevantes; a seleção ou combinação de modelos pode ser condicionada tanto à antecedência requerida quanto ao regime sazonal.

\section{Análise da Importância das Variáveis}

O gráfico de importância por ganho médio do XGBoost (Figura~\ref{fig:importancia}) mostra as variáveis que mais contribuíram para as predições do modelo. O número de interrupções do dia de origem (\texttt{interrupcoes}[$t$]) ocupa a primeira posição, refletindo que o estado imediato da rede é a informação mais preditiva para $h=1$. Na sequência aparecem \texttt{interrupcoes\_lag\_1}, \texttt{mes\_cos}, \texttt{interrupcoes\_lag\_2}, a EMA de sete dias da precipitação, \texttt{mes\_sin} e a precipitação defasada de um dia.

Esse ranking é coerente com a forte autocorrelação temporal evidenciada na análise ACF (Seção~\ref{sec:autocorrelacao}): o estado recente da série é o melhor preditor de curto prazo. As codificações harmônicas (\texttt{mes\_sin}, \texttt{mes\_cos}) indicam que a posição no ciclo sazonal do Cerrado é relevante para o modelo. As EMAs climáticas foram usadas em divisões do modelo completo, mas a ablação mostra que essa utilização interna não se converteu em ganho incremental do grupo climático no período de teste.

É importante ressaltar que o ganho mede a contribuição de cada variável para as predições do modelo, mas não um efeito causal das variáveis sobre as interrupções. Correlação temporal entre clima e falhas foi demonstrada nas seções anteriores, mas não permite inferir causalidade isolada.

\section{Papel Incremental do Clima na Previsão Diária}\label{sec:discussao_ablacao}

O achado mais incisivo da ablação é a vantagem observada da variante \textbf{histórico + calendário}: MAE 54,26 contra 61,47 do XGBoost completo. Com blocos de sete e 14 dias, o intervalo pareado permaneceu abaixo de zero; com 30 dias, o limite superior alcançou 0,11. A evidência favorece a especificação parcimoniosa nas escalas semanal e quinzenal de reamostragem, mas perde separação estrita quando se preservam segmentos mensais inteiros.

Esse resultado não contradiz a existência de associação entre condições meteorológicas e interrupções. Correlação descritiva e ganho preditivo incremental são propriedades diferentes. Chuva, vento e temperatura podem covariar com o alvo, enquanto o histórico recente das interrupções e o calendário já absorvem parte substancial da informação sazonal compartilhada. Quando atributos redundantes são acrescentados a uma amostra de 2.700 origens de treinamento, o modelo também recebe mais oportunidades de ajustar relações que não se repetem no ano de teste.

Há pelo menos quatro explicações compatíveis com o padrão observado:
\begin{enumerate}
    \item \textbf{resolução espacial}: uma única estação meteorológica não representa tempestades localizadas em todo o Distrito Federal;
    \item \textbf{agregação diária}: máximos, médias e totais diários ocultam duração, horário, trajetória e simultaneidade dos eventos;
    \item \textbf{redundância temporal}: calendário, histórico do alvo e clima compartilham sazonalidade, dificultando isolar contribuição incremental;
    \item \textbf{especificação controlada}: as variantes reutilizam hiperparâmetros selecionados para o modelo completo, sem otimização própria.
\end{enumerate}

A variante clima + calendário, sem histórico do alvo, obteve MAE 79,22 e $R^2=0{,}076$. Ela preserva algum sinal acima de uma previsão constante, mas não supera a persistência nem o modelo somente calendário. Isso mostra que as covariáveis meteorológicas disponíveis não sustentam, isoladamente, a precisão alcançada pelos modelos autorregressivos neste protocolo.

Também seria incorreto declarar histórico + calendário como novo modelo vencedor definitivo. A variante foi comparada no mesmo teste usado para produzir o achado; escolhê-la e reportá-la como modelo final usaria o teste para seleção. O procedimento correto é congelar essa hipótese e avaliá-la em dados posteriores a maio de 2025 ou em validação temporal externa. Até essa confirmação, a ablação qualifica a interpretação do modelo oficial, mas não substitui os resultados comparativos originalmente definidos.


\section{Comparação com a Literatura}

A Tabela~\ref{tab:comparacao_literatura} contextualiza os resultados deste trabalho em relação a estudos recentes da literatura internacional.

\begin{table}[H]
\caption{Comparação com Trabalhos Correlatos da Literatura}
\label{tab:comparacao_literatura}
\begin{center}
\footnotesize
\begin{tabularx}{\textwidth}{l l l X l}
\toprule
\textbf{Estudo} & \textbf{Melhor Modelo} & \textbf{Tarefa} & \textbf{Métrica Principal} & \textbf{Contexto} \\
\midrule
\textbf{Este trabalho} & XGB/LSTM/GRU & Regressão & MAE = 61,47--62,10; $R^2_{\max}=0{,}442$ & Diário, Brasil (DF) \\
\cite{zhang2023deep} & MLP & Regressão & MAE = 0,013; RMSE = 0,048 & Horário, EUA \\
\cite{smith2021hybrid} & Árvores+Vegetação & Regressão & MdAPE = 59\% & Diário, EUA \\
\cite{saha2019machine} & PSO+XGBoost & Risco de falha & Acurácia = 96,2\% & Diário, China \\
\cite{silva2021xgboost} & XGBoost/RF & Duração & Acurácia = 98,4\% & Evento, Turquia \\
\cite{roque2017weather} & AdaBoost+ & Regressão & MSE, MAE\textsuperscript{$\dagger$} & Mensal, EUA \\
\bottomrule
\multicolumn{5}{l}{\textsuperscript{$\dagger$}\scriptsize{Valores absolutos não disponíveis (paywall); autores reportam $R^2$ baixo.}} \\
\end{tabularx}
\normalsize
\end{center}
\small{Fonte: Compilado pelos autores (2026).}
\end{table}

Os $R^2$ de 0,409 (XGBoost), 0,408 (Bi-LSTM) e 0,442 (Bi-GRU) caracterizam a capacidade explicativa alcançada neste conjunto de dados, mas não permitem posicionar os modelos em uma ordenação direta de estado da arte. Os trabalhos da literatura variam em alvo, escala, resolução temporal, variáveis disponíveis, tamanho da amostra e região geográfica. Wang et al. utilizam dados socioeconômicos e de infraestrutura adicionais \cite{zhang2023deep}; Yang et al. incorporam dados de vegetação de alta resolução \cite{smith2021hybrid}; Fang et al. empregam otimização por enxame de partículas para seleção de hiperparâmetros \cite{saha2019machine}. A contribuição deste trabalho é, portanto, metodológica e aplicada: estudar o contexto do Distrito Federal e comparar modelos tabulares e redes recorrentes bidirecionais sob datas-alvo, cortes temporais e variáveis por instante controlados, explicitando também a incerteza e os limites dessa comparação.

\section{Limitações do Estudo}\label{sec:limitacoes}

O presente trabalho possui limitações que devem ser explicitadas para contextualizar adequadamente os resultados e orientar trabalhos futuros:

\begin{enumerate}
    \item \textbf{Resolução temporal diária:} A agregação diária oculta a dinâmica intra-diária das interrupções. Uma tempestade de 30 minutos no horário de pico pode causar mais impacto operacional do que chuvas distribuídas ao longo de 24 horas. Dados de resolução horária ou sub-horária permitiriam modelagens mais granulares.

    \item \textbf{Restrição geográfica:} O estudo limita-se ao Distrito Federal, cuja topografia plana e clima de Cerrado apresentam especificidades que podem não se generalizar para regiões litorâneas (expostas a ventos marítimos e salinidade) ou amazônicas (regime pluviométrico distinto).

    \item \textbf{Ausência de variáveis de vegetação:} A queda de árvores sobre a rede é uma das principais causas de interrupção, porém o estudo não incorpora dados de cobertura vegetal, índice de vegetação (NDVI) ou proximidade de árvores aos condutores. Dados de sensoriamento remoto poderiam enriquecer significativamente o modelo.

    \item \textbf{Representação espacial do vento:} A direção foi tratada com média vetorial circular e componentes seno/cosseno, evitando o artefato da transição $360^\circ/0^\circ$. A limitação remanescente é espacial: uma única estação não descreve a variação do campo de vento entre todas as regiões administrativas nem a orientação local da rede.

    \item \textbf{Estação meteorológica única:} Os dados climáticos provêm de uma única estação (A001, Brasília), que pode não representar adequadamente as condições microclimáticas de todas as regiões administrativas do DF. A incorporação de múltiplas estações ou de dados de reanálise (ERA5) proporcionaria cobertura espacial mais abrangente.

    \item \textbf{Volume de dados para Deep Learning:} As 2.687 sequências de treinamento disponíveis para as redes neurais (2.700 linhas de origem para o XGBoost), embora suficientes para modelos tabulares, são modestas para redes recorrentes profundas. A disponibilidade de séries mais longas ou de dados de múltiplas concessionárias poderia favorecer as arquiteturas neurais.

    \item \textbf{Ablação pós-hoc:} Os grupos de atributos foram comparados no mesmo período reservado para a avaliação principal. O resultado é diagnóstico e precisa ser confirmado em um novo intervalo antes de orientar seleção de modelo. Além disso, as variantes reutilizaram a configuração do XGBoost completo, sem busca específica.

    \item \textbf{Único período de teste e sensibilidade da reamostragem:} A avaliação principal cobre 365 dias consecutivos, e o \textit{block bootstrap} reutiliza segmentos desse mesmo intervalo. O procedimento quantifica a variabilidade interna preservando parte da dependência, mas não reproduz outros anos, mudanças na rede ou novos regimes meteorológicos. Além disso, alguns contrastes mudaram de classificação conforme o bloco passou de sete para 14 ou 30 dias, de modo que as conclusões inferenciais permanecem condicionadas à escala temporal adotada.
\end{enumerate}

\section{Implicações Práticas para a Neoenergia Brasília}

Os resultados deste trabalho apresentam aplicabilidade direta para a gestão operacional da concessionária:

\begin{itemize}
    \item \textbf{Sistema de alerta precoce:} O protótipo permite testar alertas para o dia seguinte, mas uma implantação deve comparar prospectivamente o XGBoost completo e a variante histórico + calendário, integrar previsão meteorológica realmente disponível na origem e calibrar limiares segundo custos de falso alarme e evento perdido. O terceiro quartil histórico pode iniciar a análise, não definir sozinho uma regra operacional.

    \item \textbf{Pré-posicionamento de equipes:} Com uma previsão para o dia seguinte, emitida após o fechamento das observações do dia corrente, a concessionária poderia apoiar o planejamento de equipes de manutenção. Para emissão intradiária ou com antecedência fixa de 24 horas, seria necessário integrar previsões meteorológicas operacionais e definir uma hora de corte operacional.

    \item \textbf{Gestão regulatória:} A capacidade de antecipar dias de maior volume pode apoiar o registro de condições meteorológicas e operacionais. A previsão, porém, não identifica a causa de cada ocorrência nem substitui a documentação técnica exigida em processos relativos a DEC e FEC.

    \item \textbf{Planejamento de investimentos:} As associações observadas para rajada, precipitação e sazonalidade podem apoiar a formulação de hipóteses para investimentos em infraestrutura resiliente. Decisões como substituição de redes aéreas ou intensificação da poda preventiva exigem validação adicional com causas das ocorrências, topologia da rede e condição dos ativos.
\end{itemize}

\section{Análise Comparativa do Custo Computacional}

Além das métricas de acurácia, a viabilidade operacional dos modelos preditivos depende de seu custo computacional, especialmente para implantação em sistemas de previsão em tempo real. A Tabela~\ref{tab:custo_computacional} sumariza o custo relativo das três arquiteturas implementadas.

\begin{table}[H]
\caption{Comparativo de Custo Computacional das Arquiteturas}
\label{tab:custo_computacional}
\begin{center}
\small
\begin{tabular}{p{4.5cm} p{3cm} p{3cm} p{3cm}}
\toprule
\textbf{Critério} & \textbf{XGBoost} & \textbf{Bi-LSTM} & \textbf{Bi-GRU} \\
\midrule
Parâmetros treináveis & $\sim$5.000 nós & 162.433 & 123.905 \\
Tempo de treino (CPU) & $<$ 2 min & $\sim$25 min & $\sim$20 min \\
Necessidade de GPU & Nenhuma & Recomendável & Recomendável \\
Tempo de inferência & $<$ 1 ms/dia & $\sim$5 ms/dia & $\sim$4 ms/dia \\
Memória necessária & $\sim$50 MB & $\sim$200 MB & $\sim$180 MB \\
\bottomrule
\end{tabular}
\normalsize
\end{center}
\small{Fonte: Estimado pelos autores com base nos experimentos realizados na Configuração 2 (Tabela~\ref{tab:hardware}), em 2026.}
\end{table}

O XGBoost apresenta vantagem operacional nos critérios medidos: treina em menos de 2 minutos em CPU convencional, não requer GPU, e realiza inferências em tempo desprezível ($<1$ ms por dia previsto). As redes recorrentes, por outro lado, demandam de 20 a 25 minutos de treinamento para 150 épocas (em CPU), com benefício de aceleração via GPU. Embora os tempos de inferência das RNNs sejam aceitáveis para previsões diárias ($\sim$5 ms), a necessidade de retrainamento periódico torna seu custo computacional observado maior no ciclo de vida do modelo.

Essa análise reforça a adequação do XGBoost para implantação operacional na Neoenergia Brasília, onde as previsões precisam ser geradas diariamente com recursos computacionais limitados e sem dependência de infraestrutura de GPU.

\section{Análise do Efeito de Amplificação pela Agregação Temporal}

Um dos achados mais consistentes deste trabalho é o \textbf{aumento numérico das correlações dos principais sinais climáticos} na transição da escala diária para a semanal e mensal (Tabela~\ref{tab:correlacoes}, Capítulo~\ref{4_Resultados}). A precipitação, por exemplo, evolui de $r = 0{,}35$ (diário) para $r = 0{,}50$ (semanal) e $r = 0{,}54$ (mensal). O comportamento não é universal: a velocidade máxima permanece próxima de zero nas três escalas. A agregação reduz variações diárias, enfatiza tendências compartilhadas e diminui o número de observações; portanto, os mecanismos físicos a seguir são hipóteses compatíveis com o padrão, não conclusões causais:

\begin{enumerate}
    \item \textbf{Possível acumulação de dano:} Chuva e vento podem estar associados a efeitos imediatos e retardados; dados sobre causa e condição dos ativos seriam necessários para verificar se esses efeitos contribuem para as correlações agregadas.

    \item \textbf{Redução da variação diária:} A agregação suaviza oscilações de curto prazo de todas as séries e pode elevar correlações ao realçar tendências comuns, sem permitir separar um componente climático de outras causas não observadas.

    \item \textbf{Possíveis efeitos de limiar:} Limiares acumulados de exposição são plausíveis, mas não foram testados neste estudo e exigem dados de ativos e causas das ocorrências.
\end{enumerate}

O padrão sugere testar modelos preditivos em \textbf{escala semanal}, ao custo de menor granularidade temporal. Para a concessionária, previsões semanais podem ser úteis ao planejamento logístico de equipes e materiais, enquanto previsões diárias permanecem mais alinhadas à operação de curto prazo.

\section{Discussão sobre a Heteroscedasticidade}

A análise do diagrama de dispersão residual (Figura~\ref{fig:scatter_hetero}, Capítulo~\ref{4_Resultados}) evidenciou uma \textbf{heteroscedasticidade} pronunciada: a variância do erro de previsão cresce com o volume de interrupções reais. Uma interpretação plausível é que dias com grande número de ocorrências reúnam condições mais heterogêneas e localizadas, não representadas integralmente pelas variáveis agregadas disponíveis. Como a base não contém a topologia da rede, o estado dos ativos ou a causa individual das ocorrências, mecanismos como falhas em cascata, envelhecimento de componentes e restrições operacionais devem ser tratados como hipóteses, e não como causas demonstradas por esta análise.

Esse padrão heteroscedástico tem duas implicações metodológicas relevantes:

\begin{enumerate}
    \item \textbf{Função de perda:} A função MSE utilizada como critério de otimização penaliza desproporcionalmente os erros em dias extremos (por ser quadrática), o que enviesa os modelos para minimizar os erros nos picos em detrimento da precisão nos dias normais. Funções de perda alternativas, como o \textit{Huber Loss} ou o \textit{Quantile Loss}, poderiam produzir modelos com distribuições residuais mais homogêneas.

    \item \textbf{Previsão probabilística:} A variância não-constante do erro torna as previsões pontuais insuficientes para a tomada de decisão operacional. Os intervalos de \textit{bootstrap} acrescentados quantificam incerteza das métricas agregadas, não a incerteza de ``250 interrupções amanhã''. Intervalos preditivos calibrados continuam sendo uma extensão necessária.
\end{enumerate}

\section{Dependência Temporal dos Resíduos}\label{sec:discussao_residuos_temporais}

O diagnóstico da Seção~\ref{sec:resultados_residuos_temporais} mostra que viés médio próximo de zero não é suficiente para caracterizar adequadamente os erros. Os intervalos de viés dos três modelos contêm zero, mas as medianas são negativas e os resíduos conservam dependência temporal. Erros positivos e negativos podem se compensar na média enquanto permanecem organizados em sequências de subestimação ou superestimação. Por isso, a ausência de viés médio claramente diferente de zero não deve ser interpretada como calibração completa.

A autocorrelação de primeira ordem foi positiva no XGBoost (0,238), na Bi-LSTM (0,304) e na Bi-GRU (0,362). Em termos práticos, parte do erro cometido em um dia tende a persistir no dia seguinte. Os testes de Ljung--Box rejeitaram a ausência conjunta de autocorrelação até 14 dias em todos os modelos ($p<10^{-7}$), indicando que nenhum deles converteu a série residual em ruído branco \cite{ljung1978measure}. O resultado é compatível com regimes operacionais e meteorológicos que duram vários dias, efeitos sazonais incompletamente representados e variáveis não observadas relacionadas à condição da rede.

Essa dependência não invalida as previsões nem identifica automaticamente a especificação correta. O teste informa que existe estrutura remanescente, mas não distingue se ela decorre de sazonalidade, mudança de nível, eventos extremos agrupados, ausência de dados espaciais ou dinâmica própria da infraestrutura. Acrescentar indiscriminadamente novas defasagens poderia apenas aumentar a complexidade e o sobreajuste. A escolha de uma correção deve ser submetida ao mesmo corte temporal e à validação fora da amostra usados na comparação principal.

Como extensões, podem ser avaliados um componente autorregressivo aplicado aos resíduos, variáveis de regime, modelos em múltiplas escalas temporais ou combinações que atualizem a previsão conforme erros recentes. Também é necessário que métodos probabilísticos preservem a ordem temporal durante a calibração, pois intervalos construídos como se os resíduos fossem independentes podem apresentar cobertura inadequada em sequências de dias críticos. Procedimentos conformais temporais e multi-horizonte constituem uma direção pertinente para tratar essa estrutura \cite{wang2024conformal}.

Para o uso operacional, a principal consequência é que o risco não se resume ao erro médio de um dia isolado. Uma sequência de subestimações pode manter equipes abaixo da necessidade durante vários dias, enquanto superestimações consecutivas podem mobilizar recursos sem retorno. O monitoramento deve acompanhar resíduos recentes, cobertura dos intervalos e desempenho por regime, acionando recalibração ou retreinamento quando houver deterioração persistente. Essa leitura conecta o diagnóstico estatístico a uma política de implantação gradual, auditável e sujeita a revisão contínua.

\section{Implicações da Tendência de Crescimento das Interrupções}

A tendência de crescimento de 65,3\% nas interrupções médias diárias ao longo de oito anos (2017--2025) constitui um achado com implicações significativas tanto para a modelagem preditiva quanto para a gestão da concessionária. Os dados disponíveis demonstram o crescimento, mas não permitem identificar suas causas. Três hipóteses podem ser investigadas em trabalhos futuros:

\begin{enumerate}
    \item \textbf{Estado e envelhecimento da infraestrutura:} Dados sobre idade, inspeção, manutenção e substituição de isoladores, condutores e transformadores permitiriam avaliar se a condição dos ativos contribui para a tendência observada.

    \item \textbf{Mudanças no regime de eventos meteorológicos:} Séries climatológicas mais longas e dados de múltiplas estações permitiriam testar se frequência, intensidade e distribuição espacial de tempestades se alteraram no período analisado.

    \item \textbf{Expansão urbana, da carga e da rede:} Informações anuais sobre unidades consumidoras, extensão da rede e carregamento dos equipamentos permitiriam verificar quanto do crescimento decorre do aumento da exposição operacional.
\end{enumerate}

Essa tendência tem uma implicação direta para os modelos preditivos: um modelo treinado com dados de 2017--2020 pode subestimar sistematicamente as interrupções em 2024--2025, uma vez que a distribuição da variável \textit{target} sofreu \textbf{deslocamento conceitual} (\textit{concept drift}). A inclusão de variáveis de tendência (codificações temporais, médias móveis de longo prazo) na engenharia de atributos mitiga parcialmente esse efeito, mas o retrainamento periódico do modelo com dados recentes permanece necessário para manter a acurácia ao longo do tempo.

\section{O El Niño 2023 como Evento Singular no Conjunto de Treinamento}

O ano de 2023, fortemente influenciado pelo fenômeno \textbf{El Niño--Oscilação Sul (ENSO)}, registrou os maiores picos de interrupções da série histórica. Como detalhado na Seção~\ref{sec:anomalia} (Capítulo~\ref{4_Resultados}), os picos de novembro--dezembro de 2023 coincidem temporalmente com episódios de precipitação intensa registrados pelo INMET. Esse período integra o conjunto de treinamento (anterior a junho de 2024); afirmações sobre o desempenho preditivo dos modelos nesse recorte não são rastreáveis a um CSV oficial de predições com essas datas, pois o conjunto de teste começa em junho de 2024.

O evento El Niño 2023--2024 foi acompanhado pelo painel conjunto de INPE, INMET, ANA e CENAD \cite{inpe2023painel}. Nesta monografia, a caracterização do período limita-se às observações meteorológicas da estação utilizada e aos registros de interrupções. Não se atribuem ao ENSO, de forma isolada, os picos observados no Distrito Federal, nem se inferem mecanismos atmosféricos ou danos específicos à infraestrutura sem dados adicionais.

O conjunto de treinamento (janeiro/2017 a maio/2024) inclui integralmente o período do El Niño 2023--2024, de modo que os modelos foram ajustados com as observações meteorológicas e elétricas registradas nesse intervalo. Contudo, o ENSO não foi incluído explicitamente como variável e o conjunto contém apenas um episódio recente dessa natureza. A generalização para futuros eventos semelhantes permanece, portanto, uma questão em aberto. Essa limitação poderia ser investigada com dados de períodos anteriores (como 2015--2016) e com a incorporação explícita de índices ENSO, como o Oceanic Niño Index (ONI).

Embora o período do El Niño 2023 integre o conjunto de treinamento (anterior ao corte em junho/2024), sua presença como evento singular de forte intensidade na série levanta uma questão relevante sobre generalização: caso futuros eventos ENSO de intensidade similar ocorram durante o período de previsão, os modelos terão aprendido com apenas um exemplar desse tipo de anomalia. A incorporação de teleconexões climáticas, como o Oceanic Niño Index (ONI), e de dados de períodos anteriores (El Niño 2015--2016) emerge como uma direção promissora para aumentar a robustez dos modelos preditivos perante perturbações climáticas de larga escala.

\section{Reflexão sobre a Escolha Metodológica: Engenharia de Atributos e Modelagem Sequencial}

Uma reflexão metodológica importante emerge da comparação entre o XGBoost e as redes recorrentes: \textbf{todos os modelos oficiais receberam as mesmas 41 variáveis por instante e compartilharam o corte temporal}. A diferença é que o XGBoost processa uma linha na origem, enquanto Bi-LSTM e Bi-GRU processam janelas sequenciais de 14 dias. A proximidade entre os modelos em $h=1$, agora acompanhada por contrastes pareados que contêm zero, mostra resultados estatisticamente inseparáveis em MAE sob essas representações distintas.

O desempenho do XGBoost na avaliação \textit{one-step} é compatível com o uso de atributos de domínio em um conjunto de dados modesto (3.066 dias após a engenharia de atributos). A ablação mostra especificamente que histórico e calendário concentram o ganho, enquanto acrescentar todo o grupo climático não melhora o teste. Redes neurais profundas possuem potencial para extrair representações hierárquicas diretamente de dados menos processados \cite{goodfellow2016deep}; essa hipótese não foi testada, pois não se realizou ablação equivalente das entradas das redes.

Essa evidência sugere que, para bases históricas limitadas, a parcimônia deve ser avaliada explicitamente. Defasagens do alvo e codificações sazonais podem oferecer mais generalização que um conjunto amplo de covariáveis correlacionadas. Variáveis meteorológicas não devem ser descartadas; devem ser melhoradas espacial e temporalmente e demonstrar ganho incremental em validação prospectiva antes de sustentar uma recomendação operacional.
